\documentclass[../../../../master.tex]{subfiles}

\begin{document}

\monographleaf{Domain 06: Reinforcement Learning \& Alignment}{Bandits, Markov Decision Processes \& Value-Based RL}{Multi-Armed Bandits, MDPs \& Dynamic Programming}{Multi-Armed \& Contextual Bandits, UCB, Thompson Sampling \& Regret Bounds}{Sutton \& Barto Ch.~2, Murphy RL Ch.~1 \& 7, Murphy Adv. Sec.~34.4}
\setlength{\emergencystretch}{2.5em}

% =============================================================================
\section{Problem Formulation \& Unified Notation}
\label{sec:bandit:formulation}
% =============================================================================

Unlike supervised learning, which receives \emph{instructive feedback} specifying the correct target independently of the learner's output, reinforcement learning relies on \emph{evaluative feedback}, where the agent observes a scalar reward $r_t$ evaluating only the selected action $a_t$ \bookcite{reinforcement_learning_introduction_richard_s_su}{Ch.~2, p.~25}. A \textbf{bandit problem} isolates this exploration--exploitation dilemma in a single-step setting where actions do not alter future environment states.

\begin{coremath}[Unified Bandit \& Contextual Bandit Formulation]
Let $\Actions = \{1, \dots, K\}$ be a finite action (arm) set. At each round $t \in \{1, \dots, T\}$:
\begin{enumerate}
  \item In a \textbf{contextual bandit}, the environment draws a context (state) $\vecb{s}_t \in \States$ from an exogenous process $p(\vecb{s}_t \mid \vecb{s}_{1:t-1})$ independent of the agent's past actions \bookcite{reinforement_learning_kevin_murphy}{Sec.~1.2.5, p.~19}; in a \textbf{context-free $K$-armed bandit}, $\States = \{\emptyset\}$ is trivial.
  \item The agent selects an action $a_t \sim \pi_t(\cdot \mid \vecb{s}_t, h_t)$ conditioned on the history $h_t = \{(\vecb{s}_\tau, a_\tau, r_\tau)\}_{\tau=1}^{t-1}$.
  \item The environment emits a stochastic reward $r_t \sim p_R(\cdot \mid \vecb{s}_t, a_t)$ with conditional expectation $R(\vecb{s}_t, a_t) \triangleq \E[r_t \mid \vecb{s}_t, a_t]$ (or $q_*(a) = \mu_a = R(a)$ when context-free).
\end{enumerate}
The finite-horizon objective is to maximize expected cumulative reward $J = \sum_{t=1}^T \E[r_t]$ (or discounted $J = \sum_{t=1}^\infty \gamma^{t-1}\E[r_t]$ with $\gamma \in (0,1)$) \bookcite{probabilisticml_kevinmurphy_advancedtopics}{Sec.~34.4.1, Eq.~34.48--34.49, pp.~1120--1121}. Letting $\pi^*(\vecb{s}) \in \argmax_{a \in \Actions} R(\vecb{s}, a)$ be the oracle policy, the \textbf{per-step regret} and \textbf{cumulative regret} are
\begin{equation}
  \ell_t \triangleq \E_{\vecb{s}_t}\!\left[R(\vecb{s}_t, \pi^*(\vecb{s}_t))\right] - \E_{\vecb{s}_t, a_t \sim \pi_t}\!\left[R(\vecb{s}_t, a_t)\right],
  \qquad
  L_T \triangleq \E\!\left[\sum_{t=1}^T \ell_t\right] = \sum_{a \in \Actions} \E[N_{T+1}(a)]\,\Delta_a,
  \label{eq:bandit:regret_decomp}
\end{equation}
where the second equality holds in the context-free case with suboptimality gap $\Delta_a \triangleq R^* - R(a)$ and pull count $N_t(a) \triangleq \sum_{\tau=1}^{t-1} \mathbb{I}(a_\tau = a)$ \bookcite{probabilisticml_kevinmurphy_advancedtopics}{Sec.~34.4.7, Eq.~34.70--34.72, pp.~1127--1128}.
\end{coremath}

\begin{intuition}[Cumulative Regret vs.\ Simple Regret]
Minimizing cumulative regret $L_T$ corresponds to ``earning while learning'': every exploratory pull of a suboptimal arm $a$ incurs immediate cost $\Delta_a$, so a consistent policy must achieve sublinear regret $L_T = o(T)$. By contrast, \textbf{pure exploration} (or \emph{best-arm identification}) minimizes only the terminal \emph{simple regret} $\ell_T$ after a fixed budget $T$, decoupling data collection from exploitation \bookcite{reinforement_learning_kevin_murphy}{Sec.~1.2.7.1, p.~22}.
\end{intuition}

\begin{crossbook}[Cross-Book Notation \& Pedagogical Reconciliation]
The three source texts adopt distinct symbols and framing conventions for the same underlying decision process. Table~\ref{tab:bandit:notation} establishes a unified correspondence.
\begin{center}
\small
\begin{tabularx}{\textwidth}{@{}l l l l X@{}}
\toprule
\textbf{Concept} & \textbf{Sutton \& Barto} & \textbf{Murphy RL} & \textbf{Murphy Adv.} & \textbf{Unified Monograph Notation} \\
\midrule
Number of arms & $k$ & $K$ or $n$ & $K$ & $K$ discrete actions $a \in \{1,\dots,K\}$ \\
Action \& reward & $A_t, R_t$ & $a_t, r_t$ & $a_t, r_t$ & $a_t \in \Actions$, $r_t \in \R$ at step $t$ \\
True arm value & $q_*(a)$ & $R(a), \mu(a)$ & $R(\vecb{s},a), \mu_a$ & $R(\vecb{s},a) = \E[r_t \mid \vecb{s}, a]$; gap $\Delta_a$ \\
Action estimate & $Q_t(a)$ & $\hat{R}_t(a), \hat{\mu}_t(a)$ & $\hat{R}_t(\vecb{s},a), \hat{\mu}_t(a)$ & $\hat{\mu}_t(a)$ (mean), $\tilde{R}_t(\vecb{s},a)$ (UCB) \\
Pull count & $N_t(a)$ & $N_t(a)$ & $N_t(a)$ & $N_t(a) = \sum_{\tau=1}^{t-1}\mathbb{I}(a_\tau=a)$ \\
Context / state & Associative cue & $o_t = w_t$ (POMDP) & $\vecb{s}_t \in \States$ & $\vecb{s}_t \in \States$ (feature vector $\vecb{x}_{t,a} \in \R^d$) \\
Total regret & (Empirical reward) & $\mathrm{Regret}_T(\pi)$ & $L_T$ & $L_T = \sum_{t=1}^T \ell_t$ \\
\bottomrule
\end{tabularx}
\captionof{table}{Unified notation across the three synthesized bandit references.}\label{tab:bandit:notation}
\end{center}
\end{crossbook}

% =============================================================================
\section{Core Mathematical Derivations \& Theorems}
\label{sec:bandit:derivations}
% =============================================================================

\subsection{Incremental Action-Value Estimation \& Nonstationary Tracking}

Storing all historical rewards requires $O(T)$ memory; incremental updates achieve $O(1)$ per-step memory and compute \bookcite{reinforcement_learning_introduction_richard_s_su}{Sec.~2.4, p.~31 (PDF p.~53)}. Given $n-1$ observations of an arm with sample average $Q_n$, the $n$-th reward $r_n$ updates the estimate via
\begin{equation}
  Q_{n+1} = \frac{1}{n}\sum_{i=1}^n r_i = Q_n + \frac{1}{n}\bigl(r_n - Q_n\bigr).
  \label{eq:bandit:incremental}
\end{equation}
Replacing $1/n$ with a constant step size $\alpha \in (0,1]$ yields an exponential recency-weighted average for nonstationary environments:
\begin{equation}
  Q_{n+1} = (1-\alpha)^n Q_1 + \sum_{i=1}^n \alpha(1-\alpha)^{n-i} r_i.
  \label{eq:bandit:exp_recency}
\end{equation}
Because $(1-\alpha)^n + \sum_{i=1}^n \alpha(1-\alpha)^{n-i} = 1$, the constant step size tracks drifting reward distributions while retaining an initial bias $(1-\alpha)^n Q_1$ \bookcite{reinforcement_learning_introduction_richard_s_su}{Sec.~2.5, p.~32}. Setting $Q_1(a)$ optimistically high (e.g., $+5$ when $q_*(a)\sim\Normal(0,1)$) induces temporary exploration under a greedy policy \bookcite{reinforcement_learning_introduction_richard_s_su}{Sec.~2.6, p.~34}. To remove initial bias while tracking nonstationary rewards, the unbiased constant-step-size trick sets $\beta_n = \alpha/\bar{o}_n$ with trace $\bar{o}_n = \bar{o}_{n-1} + \alpha(1 - \bar{o}_{n-1})$, $\bar{o}_0 = 0$ \bookcite{reinforcement_learning_introduction_richard_s_su}{Ex.~2.7, p.~35}.

\subsection{Optimism in the Face of Uncertainty: UCB1 \& Lai--Robbins Bounds}

Random $\epsilon$-greedy exploration wastes pulls on clearly inferior arms. The \textbf{Optimism in the Face of Uncertainty (OFU)} principle constructs an upper confidence bound $\tilde{R}_t(a) = \hat{\mu}_t(a) + \delta_t(a)$ such that $\tilde{R}_t(a) \ge R(a)$ with high probability, and selects $a_t = \argmax_{a \in \Actions} \tilde{R}_t(a)$ \bookcite{reinforement_learning_kevin_murphy}{Sec.~7.2.3, Eq.~7.23, p.~178}.

\begin{coremath}[Derivation of Hoeffding UCB1 \& Lai--Robbins Lower Bound]
Assume rewards for each arm $a$ are bounded in $[0,1]$ (or $1/2$-sub-Gaussian). By Hoeffding's inequality, for a fixed pull count $n = N_t(a) \ge 1$, the empirical mean $\hat{\mu}_t(a)$ satisfies
\begin{equation}
  \Prob\!\left(\mu_a - \hat{\mu}_t(a) \ge \delta\right) \le \exp\!\left(-2n\delta^2\right).
\end{equation}
Setting the failure probability to $t^{-4}$ (so that $\sum_{t=1}^\infty t^{-4} < \infty$ via a union bound over $t$) and solving $\exp(-2n\delta^2) = t^{-4}$ gives $\delta_t(a) = \sqrt{\frac{2\ln t}{N_t(a)}}$, yielding the \textbf{UCB1} rule (generalized with exploration scale $c > 0$ in \bookcite{reinforcement_learning_introduction_richard_s_su}{Eq.~2.10, p.~35} and \bookcite{probabilisticml_kevinmurphy_advancedtopics}{Eq.~34.64, p.~1125}):
\begin{equation}
  a_t = \argmax_{a \in \Actions} \left[ \hat{\mu}_t(a) + c\sqrt{\frac{\ln t}{N_t(a)}} \right], \qquad \text{(pulling any arm with } N_t(a)=0 \text{ first)}.
  \label{eq:bandit:ucb1}
\end{equation}
For any suboptimal arm ($\Delta_a > 0$), UCB1 can pull $a$ more than $\ell = \lceil 8 \ln T / \Delta_a^2 \rceil$ times only if either the optimal arm's UCB underestimates $\mu^*$ or arm $a$'s empirical mean overestimates $\mu_a + \Delta_a/2$, yielding logarithmic problem-dependent regret $L_T \le \sum_{a:\Delta_a>0} \frac{8 \ln T}{\Delta_a} + O(K)$ and minimax regret $O(\sqrt{KT\ln T})$.

This matches the fundamental \textbf{Lai--Robbins asymptotic lower bound} for any uniformly good consistent policy \bookcite{probabilisticml_kevinmurphy_advancedtopics}{Sec.~34.4.7, Eq.~34.73, p.~1128 (PDF p.~1162)}:
\begin{equation}
  \liminf_{T \to \infty} \frac{L_T}{\ln T} \ge \sum_{a:\,\Delta_a > 0} \frac{\Delta_a}{\KL{p_R(a)}{p_R(a^*)}}.
  \label{eq:bandit:lai_robbins}
\end{equation}
In the Bayesian UCB variant, one replaces the Hoeffding radius with $c\hat{\sigma}_t(a)$, where $\hat{\sigma}_t(a)$ is the posterior standard deviation (e.g., $\sqrt{\hat{\mu}_t(a)(1-\hat{\mu}_t(a))/N_t(a)}$ for Beta-Bernoulli) \bookcite{reinforement_learning_kevin_murphy}{Sec.~7.2.3.3, Eq.~7.27--7.28, p.~179}.
\end{coremath}

\subsection{Bayesian Optimal Gittins Indices \& Thompson Sampling}

In the Bayesian formulation, the agent maintains a belief state $b_t = p(\vecb{\theta} \mid h_t)$ updated deterministically via Bayes' rule, $b_t = \mathrm{BayesRule}(b_{t-1}, a_t, r_t)$ \bookcite{probabilisticml_kevinmurphy_advancedtopics}{Eq.~34.55, p.~1123}. For a context-free Beta-Bernoulli bandit with uniform prior $\mathrm{Beta}(1,1)$ and counts $(w_a, l_a)$, the finite-horizon Bayes-optimal policy satisfies the belief-state Bellman equation \bookcite{reinforement_learning_kevin_murphy}{Sec.~7.2.1.1, Eq.~7.14--7.16, p.~176}:
\begin{equation}
\begin{aligned}
  V^*(b) = \max_{a \in \Actions} \Biggl[ &\frac{w_a+1}{w_a+l_a+2}\bigl(1 + V^*(\dots, w_a+1, l_a, \dots)\bigr) \\
  &+ \left(1 - \frac{w_a+1}{w_a+l_a+2}\right) V^*(\dots, w_a, l_a+1, \dots) \Biggr].
\end{aligned}
  \label{eq:bandit:gittins_dp}
\end{equation}
While infinite-horizon discounted context-free bandits decouple via \textbf{Gittins indices}, the belief-state space scales as $O(h^{2K})$ and is provably intractable for general contextual bandits \bookcite{probabilisticml_kevinmurphy_advancedtopics}{Sec.~34.4.4, p.~1124}.

\textbf{Thompson Sampling (Probability Matching)} circumvents this intractability by choosing each action according to its posterior probability of optimality \bookcite{reinforement_learning_kevin_murphy}{Sec.~7.2.2.1, p.~177}:
\begin{equation}
  \pi_t(a \mid \vecb{s}_t, h_t) = \int \mathbb{I}\!\left(a = \argmax_{a' \in \Actions} R(\vecb{s}_t, a'; \vecb{\theta})\right) p(\vecb{\theta} \mid h_t)\, \mathrm{d}\vecb{\theta}.
  \label{eq:bandit:ts_integral}
\end{equation}
Operationally, this integral is realized without numerical integration by drawing a \emph{single} posterior sample $\tilde{\vecb{\theta}}_t \sim p(\vecb{\theta} \mid h_t)$ and acting greedily: $a_t = \argmax_{a'} R(\vecb{s}_t, a'; \tilde{\vecb{\theta}}_t)$. Despite its simplicity, Thompson Sampling achieves $O(\log T)$ problem-dependent regret and $O(\sqrt{KT\log T})$ worst-case Bayesian regret \bookcite{probabilisticml_kevinmurphy_advancedtopics}{Sec.~34.4.6--34.4.7, pp.~1126--1128}.

\monographfig[0.82\textwidth]{probabilisticml_kevinmurphy_advancedtopics_p1157_fig1.png}{Sequential Beta-Bernoulli belief-state tree for a two-armed bandit from \bookcite{probabilisticml_kevinmurphy_advancedtopics}{Fig.~34.8, p.~1123 (PDF p.~1157)}. Updating $(\alpha_a, \beta_a)$ after each pull transitions the agent deterministically through the belief-state MDP.}{fig:bandit:beta_tree}

\subsection{Contextual Bandits: LinUCB \& Off-Policy Evaluation}

When contexts $\vecb{s}_t$ carry predictive structure across actions via feature maps $\vecb{x}_{t,a} = \phi(\vecb{s}_t, a) \in \R^d$, assuming a linear reward model $R(\vecb{s}_t, a) = \vecb{\theta}_a\T \vecb{x}_{t,a} + \varepsilon_t$ ($\varepsilon_t$ $\sigma$-sub-Gaussian) enables generalization across unseen contexts \bookcite{probabilisticml_kevinmurphy_advancedtopics}{Sec.~34.4.5--34.4.6, pp.~1124--1126}.

\begin{coremath}[Ridge Posterior \& Disjoint LinUCB Decision Rule]
For each arm $a \in \Actions$ with ridge regularizer $\lambda > 0$, define the design Gram matrix and reward-weighted sufficient statistic:
\begin{equation}
  \mat{A}_a = \lambda \eye_d + \sum_{\tau < t:\, a_\tau = a} \vecb{x}_{\tau,a}\vecb{x}_{\tau,a}\T,
  \qquad
  \vecb{b}_a = \sum_{\tau < t:\, a_\tau = a} r_\tau \vecb{x}_{\tau,a},
  \qquad
  \hat{\vecb{\theta}}_a = \mat{A}_a\inv \vecb{b}_a.
  \label{eq:bandit:linucb_ridge}
\end{equation}
Under a Gaussian prior $\vecb{\theta}_a \sim \Normal(\zeros, \lambda\inv \eye_d)$ and noise variance $\sigma^2$, the posterior covariance is $\sigma^2 \mat{A}_a\inv$, so the predictive standard deviation along $\vecb{x}_{t,a}$ is proportional to $\sqrt{\vecb{x}_{t,a}\T \mat{A}_a\inv \vecb{x}_{t,a}}$. \textbf{LinUCB} selects
\begin{equation}
  a_t = \argmax_{a \in \Actions} \left[ \hat{\vecb{\theta}}_a\T \vecb{x}_{t,a} + \alpha \sqrt{\vecb{x}_{t,a}\T \mat{A}_a\inv \vecb{x}_{t,a}} \right],
  \label{eq:bandit:linucb_action}
\end{equation}
achieving regret $\tilde{O}(d\sqrt{T})$ independent of $K$ when parameters are shared across arms. Conversely, to evaluate a target policy $\pi$ offline from logged data $\{(\vecb{s}_i, a_i, r_i)\}_{i=1}^n$ collected under a stochastic logging policy $\pi_0(a_i \mid \vecb{s}_i) > 0$, one uses \textbf{Inverse Propensity Weighting (IPW)} or its variance-reduced \textbf{Doubly Robust (DR)} extension:
\begin{equation}
\begin{aligned}
  \hat{V}_{\mathrm{IPW}}(\pi) &= \frac{1}{n}\sum_{i=1}^n \frac{\mathbb{I}(\pi(\vecb{s}_i) = a_i)}{\pi_0(a_i \mid \vecb{s}_i)}\, r_i, \\
  \hat{V}_{\mathrm{DR}}(\pi) &= \frac{1}{n}\sum_{i=1}^n \left[ \hat{R}(\vecb{s}_i, \pi(\vecb{s}_i)) + \frac{\mathbb{I}(\pi(\vecb{s}_i) = a_i)}{\pi_0(a_i \mid \vecb{s}_i)}\bigl(r_i - \hat{R}(\vecb{s}_i, a_i)\bigr) \right],
\end{aligned}
  \label{eq:bandit:ipw_dr}
\end{equation}
where $\hat{V}_{\mathrm{DR}}$ is unbiased if \emph{either} the reward model $\hat{R}$ or the logging propensity $\pi_0$ is accurately specified.
\end{coremath}

\subsection{Gradient Bandits \& Baseline Variance Reduction}

Rather than estimating action values, \textbf{gradient bandit algorithms} parameterize a softmax policy over numerical preferences $H_t(a) \in \R$ \bookcite{reinforcement_learning_introduction_richard_s_su}{Sec.~2.8, Eq.~2.11--2.13, pp.~37--40 (PDF pp.~59--62)}:
\begin{equation}
  \pi_t(a) \triangleq \Prob\{a_t = a\} = \frac{e^{H_t(a)}}{\sum_{b=1}^K e^{H_t(b)}}.
  \label{eq:bandit:softmax_policy}
\end{equation}

\begin{coremath}[Exact Stochastic Gradient Ascent \& Baseline Invariance Proof]
Differentiating expected reward $\E[r_t] = \sum_{x \in \Actions} \pi_t(x) q_*(x)$ with respect to $H_t(a)$ and inserting an arbitrary action-independent scalar baseline $B_t$ (valid because $\sum_x \pd{\pi_t(x)}{H_t(a)} = \pd{}{H_t(a)} 1 = 0$) yields
\begin{equation}
  \pd{\E[r_t]}{H_t(a)}
  = \sum_{x \in \Actions} \bigl(q_*(x) - B_t\bigr) \pd{\pi_t(x)}{H_t(a)}
  = \sum_{x \in \Actions} \pi_t(x) \bigl(q_*(x) - B_t\bigr) \frac{\partial \pi_t(x)/\partial H_t(a)}{\pi_t(x)}.
\end{equation}
By the quotient rule on \eqref{eq:bandit:softmax_policy}, $\pd{\pi_t(x)}{H_t(a)} = \pi_t(x)\bigl(\mathbb{I}_{a=x} - \pi_t(a)\bigr)$. Substituting $B_t = \bar{r}_t$ (running average reward) and replacing $q_*(a_t)$ with $r_t$ (unbiased since $\E[r_t \mid a_t] = q_*(a_t)$) gives
\begin{equation}
  \pd{\E[r_t]}{H_t(a)} = \E\!\left[(r_t - \bar{r}_t)\bigl(\mathbb{I}_{a=a_t} - \pi_t(a)\bigr)\right],
  \label{eq:bandit:grad_bandit_sga}
\end{equation}
yielding the stochastic gradient ascent (REINFORCE) update $H_{t+1}(a) = H_t(a) + \alpha(r_t - \bar{r}_t)(\mathbb{I}_{a=a_t} - \pi_t(a))$ for all $a \in \Actions$.
\end{coremath}

% =============================================================================
\section{Algorithmic Implementation \& Comparative Synthesis}
\label{sec:bandit:algorithms}
% =============================================================================

Algorithm~\ref{alg:bandit:unified} unifies frequentist UCB1, Bayesian Beta-Bernoulli Thompson Sampling, and contextual LinUCB within a single online interaction loop.

\begin{algorithm}[htbp]
\caption{Unified Bandit Exploration: UCB1, Thompson Sampling \& Contextual LinUCB}
\label{alg:bandit:unified}
\KwIn{Horizon $T$, action set $\Actions=\{1,\dots,K\}$, mode $\in \{\mathrm{UCB1}, \mathrm{TS}, \mathrm{LinUCB}\}$, hyperparameters $(c, \alpha_0, \beta_0, \lambda, \alpha_{\mathrm{lin}})$}
\KwOut{Action sequence $\{a_t\}_{t=1}^T$ and updated sufficient statistics}
Initialize $N(a) \leftarrow 0$, $\hat{\mu}(a) \leftarrow 0$; $(\alpha_a, \beta_a) \leftarrow (\alpha_0, \beta_0)$; $\mat{A}_a \leftarrow \lambda \eye_d$, $\vecb{b}_a \leftarrow \zeros$ for all $a \in \Actions$\;
\For{$t = 1, 2, \dots, T$}{
  Observe context $\vecb{s}_t$ (and arm features $\vecb{x}_{t,a} \in \R^d$ if contextual)\;
  \uIf{$\mathrm{mode} = \mathrm{UCB1}$}{
    \lIf{$\exists a : N(a) = 0$}{$a_t \leftarrow a$}
    \lElse{$a_t \leftarrow \argmax_{a \in \Actions} \bigl[\hat{\mu}(a) + c\sqrt{\ln t / N(a)}\bigr]$}
  }
  \uElseIf{$\mathrm{mode} = \mathrm{TS}$ (Beta-Bernoulli)}{
    Draw $\tilde{\theta}_{t,a} \sim \mathrm{Beta}(\alpha_a, \beta_a)$ independently for all $a \in \Actions$\;
    $a_t \leftarrow \argmax_{a \in \Actions} \tilde{\theta}_{t,a}$\;
  }
  \Else(\tcp*[h]{Contextual LinUCB}){
    For all $a \in \Actions$, compute $\hat{\vecb{\theta}}_a \leftarrow \mat{A}_a\inv \vecb{b}_a$ and $u_{t,a} \leftarrow \hat{\vecb{\theta}}_a\T \vecb{x}_{t,a} + \alpha_{\mathrm{lin}} \sqrt{\vecb{x}_{t,a}\T \mat{A}_a\inv \vecb{x}_{t,a}}$\;
    $a_t \leftarrow \argmax_{a \in \Actions} u_{t,a}$\;
  }
  Execute $a_t$ and observe reward $r_t \sim p_R(\cdot \mid \vecb{s}_t, a_t)$\;
  $N(a_t) \leftarrow N(a_t) + 1$; \quad $\hat{\mu}(a_t) \leftarrow \hat{\mu}(a_t) + \frac{1}{N(a_t)}\bigl(r_t - \hat{\mu}(a_t)\bigr)$\;
  $\alpha_{a_t} \leftarrow \alpha_{a_t} + r_t$, \quad $\beta_{a_t} \leftarrow \beta_{a_t} + (1 - r_t)$\;
  $\mat{A}_{a_t} \leftarrow \mat{A}_{a_t} + \vecb{x}_{t,a_t}\vecb{x}_{t,a_t}\T$, \quad $\vecb{b}_{a_t} \leftarrow \vecb{b}_{a_t} + r_t \vecb{x}_{t,a_t}$\;
}
\end{algorithm}

\begin{center}
\small
\begin{tabularx}{\textwidth}{@{}l l l l l X@{}}
\toprule
\textbf{Algorithm} & \textbf{Per-Step Time} & \textbf{Memory} & \textbf{Regret Bound} & \textbf{Context?} & \textbf{Key Mechanism} \\
\midrule
$\epsilon$-Greedy (const.\ $\epsilon$) & $O(K)$ & $O(K)$ & $\Omega(\epsilon T)$ (linear) & No & Uniform random exploration \\
Optimistic Init. & $O(K)$ & $O(K)$ & $O(1)$ transient & No & High $Q_1(a)$ forces early trials \\
UCB1 & $O(K)$ & $O(K)$ & $O\!\left(\sum \frac{\ln T}{\Delta_a}\right)$ & No & Frequentist Hoeffding bonus \\
Thompson Sampling & $O(K)$ & $O(K)$ & $O(\sqrt{KT\ln T})$ & Yes (GP/Lin) & Posterior probability matching \\
Gradient Bandit & $O(K)$ & $O(K)$ & Asymptotic opt. & Yes (Policy) & Softmax SGA with baseline $\bar{r}_t$ \\
LinUCB & $O(Kd^2)$ & $O(Kd^2)$ & $\tilde{O}(d\sqrt{T})$ & Yes (Linear) & Ridge ellipsoid confidence set \\
\bottomrule
\end{tabularx}
\captionof{table}{Complexity and regret comparison across bandit exploration families. Using Sherman--Morrison rank-1 updates reduces LinUCB matrix inversion from $O(Kd^3)$ to $O(Kd^2)$ per step.}\label{tab:bandit:complexity}
\end{center}

\begin{crossbook}[Cross-Book Synthesis: Bayes-Optimal vs.\ Regret-Minimizing Paradigms]
Sutton \& Barto \bookcite{reinforcement_learning_introduction_richard_s_su}{Sec.~2.10, p.~42} benchmark exploration heuristics via empirical 10-armed testbed parameter studies (averaged over 2000 tasks), showing that UCB and optimistic initialization outperform fixed $\epsilon$-greedy over finite horizons. Murphy RL \bookcite{reinforement_learning_kevin_murphy}{Sec.~7.1.3, Table~7.1, p.~175} elevates this distinction to a foundational duality: \textbf{Bayes-optimal policies} (Gittins indices, BAMDPs) maximize expected return under a known prior $P_0$ via automatic posterior-driven exploration, whereas \textbf{regret-minimizing policies} (UCB, minimax algorithms) require no prior, manually inject optimism or randomization, and remain robust to adversarial or nonstationary shifts. Thompson Sampling bridges both worlds by using posterior samples to achieve minimax-optimal frequentist regret.
\end{crossbook}

% =============================================================================
\section{Visual Diagnostics \& Geometric Structure}
\label{sec:bandit:visuals}
% =============================================================================

Figure~\ref{fig:bandit:exploration_taxonomy} contrasts the four primary exploration--exploitation decision loops, highlighting how each method maps history $h_t$ (and context $\vecb{s}_t$) to action selection.

\begin{figure}[htbp]
\centering
\begin{tikzpicture}[node distance=0.9cm and 1.1cm, >=Stealth]
  \node[monoBox, text width=3.2cm] (env) {\textbf{Environment}\\Context $\vecb{s}_t$, Reward $r_t \sim p_R(\cdot \mid \vecb{s}_t, a_t)$};
  \node[monoTealBox, text width=3.0cm, below left=1.0cm and 1.8cm of env] (eps) {\textbf{$\epsilon$-Greedy}\\With prob.\ $1-\epsilon$: $\argmax_a \hat{\mu}_t(a)$; else uniform random};
  \node[monoIndigoBox, text width=3.1cm, below=1.0cm of env] (ucb) {\textbf{UCB1 (OFU)}\\$\argmax_a \hat{\mu}_t(a) + c\sqrt{\frac{\ln t}{N_t(a)}}$\\(deterministic bonus)};
  \node[monoTealBox, text width=3.1cm, below right=1.0cm and 1.8cm of env] (ts) {\textbf{Thompson Sampling}\\Draw $\tilde{\vecb{\theta}}_t \sim p(\vecb{\theta}\mid h_t)$;\\$\argmax_a R(\vecb{s}_t, a; \tilde{\vecb{\theta}}_t)$};
  \node[monoBox, text width=4.8cm, below=1.1cm of ucb] (lin) {\textbf{Contextual LinUCB}\\$\argmax_a \hat{\vecb{\theta}}_a\T \vecb{x}_{t,a} + \alpha \sqrt{\vecb{x}_{t,a}\T \mat{A}_a\inv \vecb{x}_{t,a}}$};

  \draw[monoArrow] (env.south west) -- (eps.north);
  \draw[monoArrow] (env.south) -- (ucb.north);
  \draw[monoArrow] (env.south east) -- (ts.north);
  \draw[monoArrow] (eps.south) |- (lin.west);
  \draw[monoArrow] (ucb.south) -- (lin.north);
  \draw[monoArrow] (ts.south) |- (lin.east);
\end{tikzpicture}
\caption{Taxonomy of exploration--exploitation decision rules: undirected randomization ($\epsilon$-greedy), frequentist optimism (UCB1), Bayesian probability matching (Thompson Sampling), and contextual ridge-ellipsoid optimism (LinUCB).}\label{fig:bandit:exploration_taxonomy}
\end{figure}

Figure~\ref{fig:bandit:regret_curves} illustrates the asymptotic separation between linear regret under constant $\epsilon$-greedy ($\Omega(T)$), minimax worst-case regret $O(\sqrt{KT\ln T})$, and problem-dependent logarithmic regret $O(\ln T)$ achieved by UCB1 and Thompson Sampling near the Lai--Robbins lower bound.

\begin{figure}[htbp]
\centering
\begin{tikzpicture}
\begin{axis}[
  width=0.84\textwidth,
  height=5.8cm,
  xlabel={Horizon $T$ (steps)},
  ylabel={Cumulative Regret $L_T$},
  xmin=10, xmax=5000,
  ymin=0, ymax=280,
  grid=major,
  grid style={dashed, gray!25},
  legend pos=north west,
  legend style={font=\footnotesize, fill=white, fill opacity=0.9},
  tick label style={font=\footnotesize},
  label style={font=\small}
]
  \addplot[monographCrimson, very thick, domain=10:5000, samples=100] {0.055 * x};
  \addlegendentry{Constant $\epsilon$-greedy: $\Omega(\epsilon T)$}

  \addplot[monographIndigo, very thick, dashed, domain=10:5000, samples=120] {2.6 * sqrt(x * ln(x)) / 3.2};
  \addlegendentry{Minimax bound: $O(\sqrt{KT\ln T})$}

  \addplot[monographNavy, very thick, domain=10:5000, samples=120] {14.0 * ln(x)};
  \addlegendentry{UCB1 / Thompson Sampling: $O(\ln T)$}

  \addplot[monographTeal, thick, dotted, domain=10:5000, samples=120] {9.5 * ln(x)};
  \addlegendentry{Lai--Robbins lower bound: $\sum \frac{\Delta_a}{D_{\mathrm{KL}}}\ln T$}
\end{axis}
\end{tikzpicture}
\caption{Scaling of cumulative regret $L_T$ with horizon $T$. Fixed $\epsilon$-greedy suffers linear regret because it continues pulling suboptimal arms with probability $\epsilon(K-1)/K$, whereas UCB1 and Thompson Sampling achieve logarithmic problem-dependent regret.}\label{fig:bandit:regret_curves}
\end{figure}

Empirical validation from the source texts confirms these theoretical regimes. Figure~\ref{fig:bandit:sutton_ucb} reproduces the 10-armed testbed comparison from \bookcite{reinforcement_learning_introduction_richard_s_su}{Fig.~2.4, p.~36 (PDF p.~58)}, showing UCB ($c=2$) surpassing $\epsilon$-greedy ($\epsilon=0.1$) after the initial $K=10$ forced exploration steps. Figure~\ref{fig:bandit:murphy_ctx_ts} shows linear-Gaussian contextual Thompson Sampling from \bookcite{probabilisticml_kevinmurphy_advancedtopics}{Fig.~34.10, p.~1127 (PDF p.~1161)} adapting smoothly as arm dominance reverses over time while tracking the $c\sqrt{t}$ regret envelope.

\monographfig[0.76\textwidth]{reinforcement_learning_introduction_richard_s_su_p58_fig1.png}{Average reward on the 10-armed testbed comparing UCB ($c=2$) against $\epsilon$-greedy ($\epsilon=0.1$) \bookcite{reinforcement_learning_introduction_richard_s_su}{Fig.~2.4, p.~36 (PDF p.~58)}. Notice the characteristic spike at step $K+1=11$ once every arm has been sampled once.}{fig:bandit:sutton_ucb}

\monographfig[0.88\textwidth]{probabilisticml_kevinmurphy_advancedtopics_p1161_fig1.png}{Linear-Gaussian contextual Thompson Sampling with context $\vecb{s}_t = (1, t, t^2)$ \bookcite{probabilisticml_kevinmurphy_advancedtopics}{Fig.~34.10, p.~1127 (PDF p.~1161)} (also \bookcite{reinforement_learning_kevin_murphy}{Fig.~7.1, p.~177}). (a) True nonstationary arm rewards; (b) cumulative reward showing adaptation from arm~2 to arm~0; (c) observed cumulative regret vs.\ $c\sqrt{t}$.}{fig:bandit:murphy_ctx_ts}

% =============================================================================
\section{Numerical Pitfalls, Edge Cases \& Failure Modes}
\label{sec:bandit:pitfalls}
% =============================================================================

\begin{pitfall}[Negative Feedback Loops \& Propensity Violations in Logged Bandits]
If a bandit policy exploits an inaccurate model too early, it starves alternative arms and locks into a self-reinforcing \emph{negative feedback loop} \bookcite{probabilisticml_kevinmurphy_advancedtopics}{Fig.~34.7, p.~1122 (PDF p.~1156)}. Moreover, in offline evaluation via \eqref{eq:bandit:ipw_dr}, deterministic logging policies ($\pi_0(a \mid \vecb{s}) \in \{0,1\}$) violate \emph{positivity} (common support), making IPW undefined or causing extreme variance when $\pi_0(a_i \mid \vecb{s}_i) \ll 1$. \textbf{Remedy:} Enforce stochastic logging floors $\pi_0(a \mid \vecb{s}) \ge \epsilon_{\min} > 0$, clip importance weights, or use Doubly Robust estimation.
\end{pitfall}

\begin{pitfall}[Nonstationarity Collapse in Sample Averages, Optimistic Init \& UCB1]
Sample-average updates ($\alpha_n = 1/n$) and optimistic initial values ($Q_1(a)=+5$) treat $t=1$ as privileged; if true arm values drift, sample averages adapt arbitrarily slowly and optimistic initialization exhausts its exploration drive permanently \bookcite{reinforcement_learning_introduction_richard_s_su}{Sec.~2.6, pp.~34--35 (PDF pp.~56--57)}. Similarly, standard UCB1 accumulates $N_t(a)$ forever and fails to re-explore after regime shifts. \textbf{Remedy:} Use constant step sizes with the unbiased trace $\beta_n = \alpha/\bar{o}_n$, sliding-window / discounted UCB, or Thompson Sampling with posterior forgetting.
\end{pitfall}

\begin{pitfall}[Missing Baseline in Gradient Bandits \& Singular Gram Matrices in LinUCB]
Omitting the baseline $\bar{r}_t$ in \eqref{eq:bandit:grad_bandit_sga} does not bias the expected gradient, but when rewards are shifted (e.g., $q_*(a) \approx +4$), every update pushes preferences upward, inflating variance and severely degrading convergence \bookcite{reinforcement_learning_introduction_richard_s_su}{Fig.~2.5, pp.~37--38 (PDF pp.~59--60)}. In contextual LinUCB, the unregularized Gram matrix $\sum \vecb{x}_{\tau,a}\vecb{x}_{\tau,a}\T$ is singular whenever $N_t(a) < d$. \textbf{Remedy:} Always maintain an adaptive baseline $B_t = \bar{r}_t$ in gradient bandits, and initialize $\mat{A}_a = \lambda \eye_d$ ($\lambda > 0$) with incremental Cholesky or Sherman--Morrison updates in LinUCB.
\end{pitfall}

\end{document}
